% part: intuitionistic-logic
% chapter: propositions-as-types
% section: reduction

\documentclass[../../../include/open-logic-section]{subfiles}

\begin{document}

\olfileid{int}{pty}{red}

\olsection{হ্রাস}

স্বাভাবিক নিষ্পাদনের !!{derivation}s-এ
!!a{introduction} বিধির পরেই !!a{elimination}
বিধি এলে একটি ঘুরপথ হয়, যা “হ্রাস” করা যায়।
যেমন, $\Intro{\land}$-এর পরে $\Elim{\land}$
এলে বিধিদুটি “বাতিল” হয়: $\Elim{\land}$-এর
সিদ্ধান্ত হলো \Intro{\land}-এ গড়া সংযোজনটির
দুটি সংযোজ্যের একটি, আর $\Intro{\land}$-এর
পূর্বধারণাগুলিই সেই দুটি সংযোজ্য।
তাই কোনো সংযোজ্যের !!a{derivation} দরকার
হলে সংযোজনটির ভূমিকা ও পরে নিষ্কাশন
না করেই সংশ্লিষ্ট !!{derivation} ব্যবহার করি:
\begin{gather*}
  \bottomAlignProof
  \AxiomC{}
  \RightLabel{$\delta_1$}
  \DeduceC{$!A_1$}
  \AxiomC{}
  \RightLabel{$\delta_2$}
  \DeduceC{$!A_2$}
  \RightLabel{$\Intro{\land}$}
  \BinaryInfC{$!A_1 \land !A_2$}
  \RightLabel{$\Elim{\land}_i$}
  \UnaryInfC{$!A_i$}
  \DisplayProof\bottomAlignProof
  \quad
  \redone
  \quad
  \AxiomC{}
  \RightLabel{$\delta_i$}
  \DeduceC{$!A_i$}
  \DisplayProof
\end{gather*}

এই সরলীকরণ থেকে সংশ্লিষ্ট প্রমাণপদগুলির
জন্যও অনুরূপ হ্রাস-সম্বন্ধ ভাবতে পারি।
$N_1$ যদি~$\delta_1$-এর প্রমাণপদ
এবং $N_2$ যদি~$\delta_2$-এর হয়,
তবে বাঁয়ের প্রমাণের প্রমাণপদ এক ধাপে
ডানের প্রমাণের প্রমাণপদে হ্রাস পায়:
\[
  \proj{i}{\pair{N_1}{N_2}} \redone N_i
\]
% BN-SRC-615 TARGET-CORRECTION-BEGIN
\par\smallskip\noindent\textnormal{\small\textbf{উৎসসংশোধন (BN-SRC-615):} উৎসে যুগল-নির্মাতার দুটি যুক্তি একটি বন্ধনীর ভিতরে কমা দিয়ে লেখা; \(\pair{N_1}{N_2}\)-তে বিধি ও পূর্বধারণা অনুযায়ী দুটি স্বতন্ত্র যুক্তি দেওয়া হয়েছে।}
% BN-SRC-615 TARGET-CORRECTION-END
টাইপযুক্ত ল্যাম্বডা ক্যালকুলাসে এটিই
গুণন-টাইপের \emph{বিটা-হ্রাস} বিধি।

স্বাভাবিক নিষ্পাদনে \Intro{\land}-এর
পরে \Elim{\land}-এর মতো যে বিধিযুগল
হ্রাস করা যায়, তাকে \emph{কর্তন}
বলে। টাইপযুক্ত ল্যাম্বডা ক্যালকুলাসে
$\redone$-এর বাঁয়ের পদটিকে
\emph{রিডেক্স} এবং ডানের পদটিকে
\emph{সংকুচিত রূপ} বলে। টাইপবিহীন
ল্যাম্বডা ক্যালকুলাসে শুধু
$(\lambd[x][N])Q$-কেই রিডেক্স ধরা হয়;
টাইপযুক্ত ল্যাম্বডা ক্যালকুলাসে
$\pair{N}{M}$, $\proj{i}{N}$,
$\inj[!A]{i}{N}$,
$\dcase{N}{x_1}{M_1}{x_2}{M_2}$
ও $\abort{!A}{N}$-যুক্ত পদ
যোগ হওয়ায় সংশ্লিষ্ট রিডেক্সও আছে।
% BN-SRC-616 TARGET-CORRECTION-BEGIN
\par\smallskip\noindent\textnormal{\small\textbf{উৎসসংশোধন (BN-SRC-616):} অন্তর্ভুক্তি-নির্মাতার উৎসরূপ \(\inj{i}{!A}{N}\) পাশের tN2i বিধির ঐচ্ছিক টাইপ-যুক্তির রীতি মানে না; এখানে \(\inj[!A]{i}{N}\) করা হয়েছে।}
% BN-SRC-616 TARGET-CORRECTION-END

অনুরূপভাবে বিযোজনের হ্রাস:
\begin{gather*}
  \bottomAlignProof
  \AxiomC{}
  \RightLabel{$\delta$}
  \DeduceC{$!A_i$}
  \RightLabel{$\Intro{\lor}$}
  \UnaryInfC{$!A_1 \lor !A_2$}
  \AxiomC{$\Discharge{!A_1}{x_1}$}
  \RightLabel{$\delta_1$}
  \DeduceC{$!C$}
  \AxiomC{$\Discharge{!A_2}{x_2}$}
  \RightLabel{$\delta_2$}
  \DeduceC{$!C$}
  \DischargeRule{$\Elim{\lor}$}{x_1,x_2}
  \TrinaryInfC{$!C$}
  \DisplayProof\bottomAlignProof
  \quad
  \redone
  \quad
  \AxiomC{}
  \RightLabel{$\delta$}
  \DeduceC{$!A_i$}
  \RightLabel{$\delta_i$}
  \DeduceC{$!C$}
  \DisplayProof
\end{gather*}
এর অনুরূপ প্রমাণপদের হ্রাস:
\begin{gather*}
  \dcase{\inj[!A_i]{i}{M}}{x_1}{N_1}{x_2}{N_2} \redone \Subst{N_i}{M}{x_i}
\end{gather*}
এখানে $M$ হলো~$\delta$-র প্রমাণপদ,
আর $N_i$ হলো~$\delta_i$-র প্রমাণপদ।
সমষ্টি-টাইপের বিটা-হ্রাস বিধি এটিই।
এখানে $\Subst{M}{N_i}{x}$ মানে
$M$-এ চলরাশি~$x$-এর সব মুক্ত
আবির্ভাবের জায়গায়~$N_i$ বসানো।

অপেক্ষক-টাইপের হ্রাস $\Intro\lif$-এর
পরে $\Elim\lif$-এর ঘুরপথ সরানোর
অনুরূপ:
\begin{gather*}
  \bottomAlignProof
  \AxiomC{$\Discharge{!A}{x}$}
  \RightLabel{$\delta$}
  \DeduceC{$!B$}
  \DischargeRule{$\Intro{\lif}$}{x}
  \UnaryInfC{$!A \lif !B$}
  \AxiomC{}
  \RightLabel{$\delta'$}
  \DeduceC{$!A$}
  \RightLabel{$\Elim{\lif}$}
  \BinaryInfC{$!B$}
  \DisplayProof\bottomAlignProof
  \quad
  \redone
  \quad
  \AxiomC{}
  \RightLabel{$\delta'$}
  \DeduceC{$!A$}
  \RightLabel{$\delta$}
  \DeduceC{$!B$}
  \DisplayProof
\end{gather*}
প্রমাণপদে এটি সাধারণ বিটা-হ্রাস:
\begin{gather*}
  (\lambd[\typeof{x}{!A}][N]M)
  \redone \Subst{N}{M}{x}
\end{gather*}

!!{proof}s-এর হ্রাস-রূপান্তরের
পাশাপাশি বিন্যাস-রূপান্তরও
বিবেচনা করি। কোনো (উপ-)!!{proof}s
যদি \Elim{\lor}-এর পরে আরেকটি
!!{elimination} বিধিতে শেষ হয়,
তাতে এই রূপান্তর প্রযোজ্য; যেমন:
\begin{multline*}
  \bottomAlignProof
  \AxiomC{}
  \RightLabel{$\delta$}
  \DeduceC{$!A_i$}
  \RightLabel{$\Intro{\lor}$}
  \UnaryInfC{$!A_1 \lor !A_2$}
  \AxiomC{$\Discharge{!A_1}{x_1}$}
  \RightLabel{$\delta_1$}
  \DeduceC{$!C_1 \land !C_2$}
  \AxiomC{$\Discharge{!A_2}{x_2}$}
  \RightLabel{$\delta_2$}
  \DeduceC{$!C_1 \land !C_2$}
  \DischargeRule{$\Elim{\lor}$}{x_1,x_2}
  \TrinaryInfC{$!C_1 \land !C_2$}
  \RightLabel{\Elim{\land}[i]}
  \UnaryInfC{$!C_i$}
  \DisplayProof\bottomAlignProof
  \quad
  \redone
  \\
  \AxiomC{}
  \RightLabel{$\delta$}
  \DeduceC{$!A_i$}
  \RightLabel{$\Intro{\lor}$}
  \UnaryInfC{$!A_1 \lor !A_2$}
  \AxiomC{$\Discharge{!A_1}{x_1}$}
  \RightLabel{$\delta_1$}
  \DeduceC{$!C_1 \land !C_2$}
  \RightLabel{\Elim{\land}[i]}
  \UnaryInfC{$!C_i$}
  \AxiomC{$\Discharge{!A_2}{x_2}$}
  \RightLabel{$\delta_2$}
  \DeduceC{$!C_1 \land !C_2$}
  \RightLabel{\Elim{\land}[i]}
  \UnaryInfC{$!C_i$}
  \DischargeRule{$\Elim{\lor}$}{x_1,x_2}
  \TrinaryInfC{$!C_i$}
  \DisplayProof
\end{multline*}
$\delta$, $\delta_1$ ও $\delta_2$-র
প্রমাণপদ যথাক্রমে $M$, $N_1$ ও
$N_2$ হলে প্রমাণপদ, অর্থাৎ
$\lambd$-পদের অনুরূপ হ্রাসবিধি:
\[
  \proj{i}{\dcase{M}{x_1}{N_1}{x_2}{N_2}} \redone
\dcase{M}{x_1}{\proj{i}{N_1}}{x_2}{\proj{i}{N_2}}
\]

!!{proof}s-এর শুধু শেষেই নয়,
!!{proof}s-এর ভিতরেও ঘুরপথ সরানো যায়:
ঘুরপথে শেষ হওয়া কোনো উপ-!!{proof}-এ
হ্রাস-রূপান্তর প্রয়োগ করি।
অনুরূপভাবে $\lambd$-পদের
উপপদেও $\beta$-হ্রাস প্রযোজ্য।
$N$ যদি $M$-এর উপপদ হয় এবং
$N \redone N'$ হয়, তবে
$M$-এ উপপদ~$N$-এর জায়গায়
$N'$ বসিয়ে $M'$ পাই। তখন
$M \redone M'$-ও লিখি।
$M = M_1 \redone M_2
\redone M_3 \redone \dots
\redone M_k = M'$ রূপের
অনুক্রম থাকলে (এমনকি $k=1$,
অর্থাৎ $M=M'$ হলেও)
$M \red M'$ লিখি।
% BN-SRC-617 TARGET-CORRECTION-BEGIN
\par\smallskip\noindent\textnormal{\small\textbf{উৎসসংশোধন (BN-SRC-617):} হ্রাস-অনুক্রমে উৎসের \(M_2\) পরপর দুবার ছিল; তৃতীয় পদটি \(M_3\) করা হয়েছে। শূন্য-ধাপের \(k=1\) ক্ষেত্র অক্ষত।}
% BN-SRC-617 TARGET-CORRECTION-END

যে !!{proof}-এর কোনো
উপ-!!{proof}-এ কোনো রূপান্তর
প্রয়োগ করা যায় না, তাকে
\emph{স্বাভাবিক} বলে।
অনুরূপভাবে যে $\lambd$-পদের
কোনো উপপদে রূপান্তর প্রযোজ্য
নয়, সেটিও \emph{স্বাভাবিক}।
স্বাভাবিক $\lambd$-পদকে
\emph{স্বাভাবিক রূপ}-ও বলে।

\end{document}
